\pdfminorversion=4
\documentclass[xcolor=dvipsnames]{beamer} 
%\usetheme{Warsaw} 
\usecolortheme[RGB={0,99,0}]{structure}
\usetheme{Darmstadt} 
\usepackage[utf8]{inputenc}
\usepackage[spanish]{babel}
%\providecommand{\abs}[1]{\lvert#1\rvert}
%\renewcommand{\frametitle}{\frametitle{\subsecname}} 
%\setbeamercovered{transparent} 
\usepackage{tikz}
%\usepackage{enumitem}
%\setenumerate[1]{label=\arabic*.}
\newcounter{ResumeEnumerate}
\begin{document}
\title{ScipyCon Argentina 2013.}
\logo{  \includegraphics[width=1cm,height=1cm,keepaspectratio]{Imagenes/logo.png}}
\subtitle{Conferencia de Python en la Ciencia - Puerto Madryn, del 16 al 18 de mayo}
\date{\scalebox{3}{\insertlogo}}

%---------------------------------------------------------------------------
\begin{frame}
\titlepage
\end{frame}

%---------------------------------------------------------------------------
\section{Python es \ldots}

\begin{frame}{¿Qu\'e es Python ?}
	Python es un lenguaje de programación dinámico de proposito general. Algunas de sus características son:
	\begin{itemize}
		\item Fácil de aprender (es utilizado para aprender a programar desde la UBA hasta el MIT).
		\item Es todo terreno. 
		\item Es libre, con un ecosistema libre.
		\item Interactivo, pueden verse resultados rápidamente.
		\item Puede comunicarse con otros lenguajes fácilmente.
		\item Argentina tiene la comunidad hispanohablante más grande.
		\item Viene instalado en las máquinas de \emph{conectar igualdad}.
	\end{itemize}
% \begin{figure}[h]
%	\begin{center}
%		\includegraphics[scale = 0.4]{venn.png}
%	\end{center}
	%\caption{Landmarks de tipo II, imagen obtenida de ~\cite{Gonzalez2011} }
%\end{figure}
%Hacer diagrama de venn
\end{frame}

%---------------------------------------------------------------------------
\begin{frame}{¿Quien usa Python?}

 \begin{figure}[h]
%	\begin{center}
		\includegraphics[scale = 0.25]{Imagenes/quienusa.png}
%	\end{center}
	%\caption{Landmarks de tipo II, imagen obtenida de ~\cite{Gonzalez2011} }
\end{figure}
%
\end{frame}

%---------------------------------------------------------------------------
\begin{frame}{Qu\'e podemos hacer con Python?}

 \begin{figure}[h]
%	\begin{center}
		\includegraphics[scale = 0.25]{Imagenes/quehacer2.png}
%	\end{center}
	%\caption{Landmarks de tipo II, imagen obtenida de ~\cite{Gonzalez2011} }
\end{figure}
\end{frame}

\begin{frame}{Qu\'e podemos hacer con Python? (Cont.)}

 \begin{figure}[h]
%	\begin{center}
		\includegraphics[scale = 0.25]{Imagenes/quehacer4.png}
%	\end{center}
	%\caption{Landmarks de tipo II, imagen obtenida de ~\cite{Gonzalez2011} }
\end{figure}
\end{frame}



%---------------------------------------------------------------------------
\begin{frame}{Qu\'e podemos hacer con Python? (Cont.)}

 \begin{figure}[h]
%	\begin{center}
		\includegraphics[scale = 0.25]{Imagenes/quehacer3.png}
%	\end{center}
	%\caption{Landmarks de tipo II, imagen obtenida de ~\cite{Gonzalez2011} }
\end{figure}
\end{frame}


%---------------------------------------------------------------------------
\section{ScipyCon es \ldots}
\begin{frame}{Qu\'e es SciPyCon Argentina?}
Con esta conferencia se busca:
	\begin{itemize}
	\item Intercambiar las experiencias y dar a conocer las ventajas del lenguaje Python en la comunidad científica argentina.
	\item Presentar las herramientas y bibliotecas de uso científico desarrolladas por la comunidad mundial.
	\item Introducir los aspectos de Python que articulan la educación, la ingeniería y la ciencia.
	\item Generar un marco apropiado para concretar encuentros, foros de discusión y proyectos educativos y de investigación y desarrollo científico relacionados con Python.
\end{itemize}
\end{frame}

%---------------------------------------------------------------------------
\begin{frame}{Qu\'e vamos a ver en la SciPyCon?}
 \begin{figure}[h]
%	\begin{center}
		\includegraphics[scale = 0.25]{Imagenes/quehacer.png}
%	\end{center}
	%\caption{Landmarks de tipo II, imagen obtenida de ~\cite{Gonzalez2011} }
\end{figure}
%
\end{frame}

%---------------------------------------------------------------------------
\begin{frame}{Cómo esta compuesta la SciPyCon Argentina?}
	\begin{itemize}
		\item Talleres (AI e Introducción a numpy + scipy + ipython).
		\item Charlas (Biopython, seguridad, sistemas SCADA, etc).
		\item Plenarias(dadas por Invitados Internacionales).
		\item Sesión de posters con pequeñas charlas.
	\end{itemize}
\end{frame}
%---------------------------------------------------------------------------
\begin{frame}{Quíen participa de la SciPyCon?}
	\begin{itemize}
		\item Invitado internacional Prabhu Ramachandran (Enthought Inc., India).
		\item Científicos de todas las ramas, biólogos, ingenieros, estadístas, matemáticos, informáticos, físicos.
		\item Docentes y alumnos universitarios, terciarios y de nivel medio.
		\item Público en general. 
	\end{itemize}
\end{frame}

%---------------------------------------------------------------------------
\section{Contacto y Preguntas ...}
\begin{frame}{Contacto}
	\begin{itemize}
		\item Inscripciones y Presentación de trabajos: \url{www.scipycon.com.ar}.
		\item Consultas: scipyconar2013@gmail.com
	\end{itemize}
	\begin{center}
	      \huge{Muchas Gracias y los esperamos!}
        \end{center}
 \begin{figure}[h]
%	\begin{center}
		\includegraphics[scale = 0.25]{Imagenes/org.png}
%	\end{center}
	%\caption{Landmarks de tipo II, imagen obtenida de ~\cite{Gonzalez2011} }
\end{figure}
\end{frame}

\end{document}

